% ===== PRÉAMBULE CANONIQUE — à coller VERBATIM en tête de CHAQUE .tex =====
\documentclass[11pt,a4paper]{article}
\usepackage[utf8]{inputenc}\usepackage[T1]{fontenc}\usepackage{lmodern}
\usepackage[french]{babel}
\usepackage{amsmath,amssymb,mathtools}\usepackage{icomma}
\usepackage{siunitx}\sisetup{locale=FR}  % \qty{}{}, \si{}, \ang{}
\usepackage{xcolor}\usepackage{colortbl}
\usepackage{tikz}\usepackage{pgfplots}\pgfplotsset{compat=1.18}
\usepackage[siunitx,europeanresistors]{circuitikz} % circuits (élec.)
\usepackage[most]{tcolorbox}\usepackage{enumitem}\usepackage{array,booktabs}
\usepackage[a4paper,margin=2.2cm]{geometry}
\usetikzlibrary{arrows.meta,calc,angles,quotes,positioning,patterns,%
  backgrounds,shapes.geometric,decorations.pathmorphing,decorations.markings,intersections}
\definecolor{primary}{RGB}{222,49,99}\definecolor{primarydark}{RGB}{150,22,55}
\definecolor{defcol}{RGB}{0,114,178}\definecolor{methcol}{RGB}{52,92,148}
\definecolor{excol}{RGB}{0,135,90}\definecolor{attncol}{RGB}{200,40,55}
\tcbset{boxbase/.style={enhanced,breakable,boxrule=0.9pt,arc=2.5pt,
  left=3mm,right=3mm,top=2.3mm,bottom=2.3mm,fonttitle=\bfseries,coltitle=white}}
% --- boîtes TUTORIEL (noms EXACTS, ne pas renommer) ---
\newtcolorbox{cle}[1][]{boxbase,colback=defcol!6,colframe=defcol,
  title={\textsc{À retenir}\ifx\relax#1\relax\relax\else\textmd{ : \itshape #1}\fi}}
\newtcolorbox{methode}[1][]{boxbase,colback=methcol!7,colframe=methcol,sharp corners,
  title={\textsc{Méthode}\ifx\relax#1\relax\relax\else\textmd{ --- \itshape #1}\fi}}
\newtcolorbox{exemplebox}[1][]{boxbase,colback=excol!5,colframe=excol,
  title={\textsc{Exemple}\ifx\relax#1\relax\relax\else\textmd{ : \itshape #1}\fi}}
\newtcolorbox{piege}[1][]{boxbase,colback=attncol!6,colframe=attncol,
  title={\textsc{Piège}\ifx\relax#1\relax\relax\else\textmd{ : \itshape #1}\fi}}
\newtcolorbox{remarque}[1][]{boxbase,colback=black!4,colframe=black!45,
  title={\textsc{Remarque}\ifx\relax#1\relax\relax\else\textmd{ : \itshape #1}\fi}}
% --- boîtes EXERCICE / CORRIGÉ (noms EXACTS, auto-comptées) ---
\newtcolorbox[auto counter]{exo}[1][]{boxbase,colback=primary!4,colframe=primary,
  title={\textsc{Exercice}~\thetcbcounter\ifx\relax#1\relax\relax\else\textmd{ --- \itshape #1}\fi}}
\newtcolorbox[auto counter]{corrige}[1][]{boxbase,colback=excol!4,colframe=excol,
  title={\textsc{Corrigé}~\thetcbcounter\ifx\relax#1\relax\relax\else\textmd{ --- \itshape #1}\fi}}
\newcommand{\vv}[1]{\vec{#1}}\newcommand{\ds}{\displaystyle}
\newcommand{\R}{\mathbb{R}}\newcommand{\N}{\mathbb{N}}\newcommand{\Z}{\mathbb{Z}}
% =========================================================================
\begin{document}

\begin{corrige}[le rebond de base]
Loi de la réflexion : $\hat{r}=\hat{\imath}$, les deux angles étant mesurés par rapport à la normale.
\begin{enumerate}[label=\alph*)]
  \item $\hat{r}=\hat{\imath}=\ang{28}$.
  \item L'angle entre le rayon incident et le rayon réfléchi vaut $\hat{\imath}+\hat{r}=2\hat{\imath}
        =2\times\ang{28}=\ang{56}$.
\end{enumerate}
\end{corrige}

\begin{corrige}[angle donné par rapport au miroir]
La normale est perpendiculaire au miroir : il faut convertir l'angle donné.
\begin{enumerate}[label=\alph*)]
  \item $\hat{\imath}=\ang{90}-\ang{35}=\ang{55}$ (angle par rapport à la normale).
  \item $\hat{r}=\hat{\imath}=\ang{55}$. \emph{Piège évité} : prendre directement $\ang{35}$ aurait
        été faux.
\end{enumerate}
\end{corrige}

\begin{corrige}[miroir posé au sol]
La normale est verticale (perpendiculaire au sol horizontal).
\begin{enumerate}[label=\alph*)]
  \item $\hat{r}=\hat{\imath}=\ang{60}$.
  \item Le rayon réfléchi fait avec la normale un angle de $\ang{60}$, donc avec le sol il fait
        l'angle complémentaire : $\ang{90}-\ang{60}=\ang{30}$.
\end{enumerate}
\end{corrige}

\begin{corrige}[incidence perpendiculaire]
Ici le rayon arrive le long de la normale, donc $\hat{\imath}=\ang{0}$.
\begin{enumerate}[label=\alph*)]
  \item $\hat{r}=\hat{\imath}=\ang{0}$.
  \item Le rayon réfléchi repart \textbf{sur lui-même}, exactement en sens inverse du rayon incident
        (demi-tour). C'est le seul cas où le rayon revient sur sa propre trajectoire.
\end{enumerate}
\end{corrige}

\begin{corrige}[rayon réfléchi perpendiculaire à l'incident]
L'angle entre le rayon incident et le rayon réfléchi vaut $2\hat{\imath}$. On veut qu'il soit égal
à $\ang{90}$ :
\[ 2\hat{\imath}=\ang{90}\;\Longrightarrow\; \hat{\imath}=\ang{45}. \]
Un rayon frappant le miroir à $\ang{45}$ de la normale repart donc à angle droit de sa direction
d'arrivée — principe utilisé pour renvoyer un faisceau « en équerre ».
\end{corrige}

\end{document}
